\documentclass[10pt]{article}
\usepackage{cornell-notes}
\title{Chapter 10: Digital Logic}\author{}\date{}
\setDocTitle{Chapter 10: Digital Logic}\setDocSubtitle{Electronics}\setDocOwner{Electronics Cornell Notes}\setDocDate{\today}
\setCornellCollection{Electronics Cornell Notes}\setCornellUnitType{Chapter}\setCornellUnitNumber{10}\setCornellUnitTitle{Digital Logic}
\hypersetup{pdftitle={Chapter 10: Digital Logic},pdfsubject={Cornell notes},pdfauthor={}}
\begin{document}\maketitle
\begin{cornell}\noterow{Core idea}{Digital logic depends on voltage thresholds, timing margins, Boolean behavior, and noise immunity.}\noterow{Validation}{Check levels, setup and hold timing, fanout, and loading.}\end{cornell}
\begin{summarybox}{Section summary}Logic values are electrical ranges with time-dependent constraints.\end{summarybox}
\end{document}
